%======================================================================
%  Title page, in the same layout as the "Nonlinear Problems in
%  Mechanics" notes: two gray rules around a large sans-serif bold
%  title, author and laboratory below, a small emblem of the running
%  example of the book, and the date.
%
%  The emblem is the quarter plate with a hole of Chapter 2, with the
%  symmetry conditions in blue (kinematic data) and the loaded edge in
%  green (static data) -- the same colour convention as the Tonti
%  diagrams of the companion notes.
%======================================================================
\begin{titlepage}
\sffamily\setlength{\parindent}{0pt}
\null\vspace{3.2cm}
{\color{gray!70!black}\rule{\textwidth}{0.8pt}}\par\vspace{18pt}
{\fontsize{30}{36}\selectfont\bfseries Introduction to \texttt{Cast3M}\\[4pt]
 programming\par}
\vspace{14pt}
{\LARGE Lecture notes\par}
\vspace{18pt}
{\color{gray!70!black}\rule{\textwidth}{0.8pt}}\par
\vspace{2.4cm}
{\Large Andrei Constantinescu\par}
\vspace{6pt}
{\normalsize\texttt{andrei.constantinescu@polytechnique.edu}\par}
\vspace{10pt}
{\normalsize Laboratoire de M\'ecanique des Solides\\
CNRS UMR 7649, \'Ecole Polytechnique\\
Institut Polytechnique de Paris, 91120 Palaiseau\par}
\vfill
\begin{center}
\begin{tikzpicture}[scale=0.75,>=Stealth]
  % the quarter plate with a hole: the running example of Chapter 2
  \begin{scope}
    \clip (0,0) rectangle (4,2.6);
    \fill[gray!10] (0,0) rectangle (4,2.6);
    \fill[white] (0,0) circle (1.1);
    \draw[gray!45,thin] (0,0) circle (1.1);
    % a coarse structured mesh, radial and circumferential
    \foreach \a in {0,15,...,90}
      \draw[gray!45,thin] (\a:1.1) -- (\a:4.8);
    \foreach \r in {1.6,2.2,2.8,3.4,4.0}
      \draw[gray!45,thin] (0:\r) arc (0:90:\r);
  \end{scope}
  % the free top edge, and the hole
  \draw[gray!55,thin] (0,2.6) -- (4,2.6);
  % symmetry conditions (kinematic data) in blue
  \draw[line width=1.6pt,cu] (1.1,0) -- (4,0);
  \draw[line width=1.6pt,cu] (0,1.1) -- (0,2.6);
  % loaded edge (static data) in green
  \draw[line width=1.6pt,cphi] (4,0) -- (4,2.6);
  \foreach \y in {0.25,0.65,...,2.55}
    \draw[->,cphi] (4.15,\y) -- (4.75,\y);
\end{tikzpicture}
\end{center}
\vspace{0.9cm}
\begin{small}
\textbf{Work in progress.} These notes are being written and rewritten as the
lectures they accompany evolve, and they are certainly not free of errors --
in the text, in the equations, and in the \textsc{gibiane} examples. Please
send corrections, comments and suggestions to the author at the address
above: they are genuinely welcome, and they are what makes the next version
better than this one.
\end{small}
\par\vspace{10pt}
{\small Draft -- \today\par}
\end{titlepage}
\cleardoublepage
